\input{../tex-inputs/latex-leseansicht-vorspann}

\section[Gustav Schwarzkopf an Arthur Schnitzler, 10. 1. 1909]{L04317 Gustav Schwarzkopf an Arthur Schnitzler, 10. 1. 1909}
\nopagebreak\mylabel{L04317v}
\rehead{ }\normalsize\beginnumbering\briefempfaengerindex{Schnitzler, Arthur@\textsc{Schnitzler, Arthur}!zzzSchwarzkopf, Gustav@\emph{von Gustav Schwarzkopf}!1909-01-101@{10. 1. 1909}|(be}
\toendnotes[C]{\smallbreak\pagebreak[2]}
\correspDesc{Versand  durch Gustav Schwarzkopf am 10. 1. 1909 in Wien
\newline{}Übermittlung  am 11. 1. 1909 in Wien
\newline{}Erhalt  durch Arthur Schnitzler im Zeitraum [10. 1. 1909 – 13. 1. 1909?] in Wien}\toendnotes[C]{\smallbreak}
\Standort{CUL, Schnitzler, B 96a.}
\physDesc{Postkarte, 732 Zeichen
\newline{}Handschrift: schwarze Tinte, deutsche Kurrent
\newline{}Versand: Stempel: »\nobreak{}\oindex{I., Innere Stadt@\textbf{I., Innere Stadt}, \emph{Verwaltungsgebiet}|pwk}1/\textsubscript{1} Wien 8, 11. I. 09, IX\nobreak{}«.  
\newline{}Schnitzler: mit Bleistift Vermerk: »\textsc{Schwarzkopf}« }\toendnotes[C]{\smallbreak}\pstart{}{\pb}I. Tiefer Graben 17\oindex{Wien@\textbf{Wien}!I., Innere Stadt@\textbf{I., Innere Stadt}!Tiefer Graben 17@\textbf{Tiefer Graben 17}, \emph{Wohngebäude}|pw}.\pend{}\pstart{}II St. Th. 6.\pend{}{\bigskip}\pstart{}Herrn\pend{}\pstart{}\textsc{D\textsuperscript{r} Arthur Schnitzler}\pend{}\pstart{}\textsc{Wien\oindex{Wien@\textbf{Wien}, \emph{Verwaltungsgebiet}|pw}}\pend{}\pstart{}XVIII Spöttelgaſſe 7\oindex{Wien@\textbf{Wien}!XVIII., Währing@\textbf{XVIII., Währing}!Edmund-Weiß-Gasse 7@\textbf{Edmund-Weiß-Gasse 7}, \emph{Wohngebäude}|pw}.\pend{}{\bigskip}\vspace{1em}
\pstart
           \raggedleft{}{\pb}10. 1. 09\pend
           \vspace{0.5em}
\pstart
           Lieber Arthur, ich danke Ihnen beſtens für die beiden Sitze\eventindex{Johann Strauß-Theater@\textbf{Johann Strauß-Theater}!Concordia-Matinee, 10.1.1909@Concordia-Matinee, 10.1.1909|pwv}. Es war eine recht anregende Vorſtellung\eventindex{Johann Strauß-Theater@\textbf{Johann Strauß-Theater}!Concordia-Matinee, 10.1.1909@Concordia-Matinee, 10.1.1909|pwv} und bei »Anatoles«
                  Hochzeitsmorgen\pwindex{Schnitzler, Arthur 15.\,5.\,1862 Wien – 21.\,10.\,1931 ebd.@\textsc{Schnitzler, Arthur} (15.\,5.\,1862 Wien – 21.\,10.\,1931 ebd.), \emph{Schriftsteller, Mediziner}!Anatols Hochzeitsmorgen@\strich\emph{Anatols Hochzeitsmorgen}|pw} wurde wirklich viel gelacht. Schon der Umſtand, daß es Morgen
               und licht war, während die anderen Stücke\pwindex{\textcolor{red}{\textsuperscript{XXXX indx1}}!Besuch in der Dämmerung@\strich\emph{Besuch in der Dämmerung}|pwv}\pwindex{\textcolor{red}{\textsuperscript{XXXX indx1}}!Eine florentinische Tragödie@\strich\emph{Eine florentinische Tragödie}|pwv}\pwindex{\textcolor{red}{\textsuperscript{XXXX indx1}}!Pechvogel@\strich\emph{Der Pechvogel}|pwv} alle bei ganz
               verdunkelter Bühne geſpielt wurden, wirkte erlöſend, Da{\geminationn}
               war es auch erfreulich, daß in dem Stückchen\pwindex{Schnitzler, Arthur 15.\,5.\,1862 Wien – 21.\,10.\,1931 ebd.@\textsc{Schnitzler, Arthur} (15.\,5.\,1862 Wien – 21.\,10.\,1931 ebd.), \emph{Schriftsteller, Mediziner}!Anatols Hochzeitsmorgen@\strich\emph{Anatols Hochzeitsmorgen}|pwv}{ }ſich nicht
               ein einziger kleiner Todesfall ereignet, während die andern drei Stücke\pwindex{\textcolor{red}{\textsuperscript{XXXX indx1}}!Besuch in der Dämmerung@\strich\emph{Besuch in der Dämmerung}|pw}\pwindex{\textcolor{red}{\textsuperscript{XXXX indx1}}!Eine florentinische Tragödie@\strich\emph{Eine florentinische Tragödie}|pw}\pwindex{\textcolor{red}{\textsuperscript{XXXX indx1}}!Pechvogel@\strich\emph{Der Pechvogel}|pw} – mäßig berechnet – 5 Tote verbrauchen, was für eine
               Nachmittagsvorſtellung entſchieden zu viel sind für die Verdauungen gleich nach dem
               Mittageſſen u nicht{\pb}zuträglich
               ist. –\pend
           
\pstart
           Ihre Frau\pwindex{Schnitzler, Olga 17.\,1.\,1882 Wien – 13.\,1.\,1970 Lugano@\textsc{Schnitzler, Olga} (17.\,1.\,1882 Wien – 13.\,1.\,1970 Lugano), \emph{Schauspielerin, Sängerin}|pwv} und{\\[\baselineskip]}Sie
               herzlichſt grüßend{\\[\baselineskip]}Ihr{\\[\baselineskip]}\spacefill\mbox{Gustav}\pend
           \leftskip=0em{}\selectlanguage{ngerman}\endnumbering\briefempfaengerindex{Schnitzler, Arthur@\textsc{Schnitzler, Arthur}!zzzSchwarzkopf, Gustav@\emph{von Gustav Schwarzkopf}!1909-01-101@{10. 1. 1909}|)be}\mylabel{L04317h}
\begin{anhang}
\end{anhang}\newcommand{\dateiname}{L04317}\newcommand{\titel}{Gustav Schwarzkopf an Arthur Schnitzler, 10. 1. 1909}\newcommand{\editorInnen}{Herausgegeben von Jahnke, SelmaMüller, Martin Anton}\input{../tex-inputs/latex-leseansicht-abspann}
